% Generated from validated Article Draft Result. Do not hand-edit; revise the structured draft instead.

\subsection{training・adaptation・inference・servingは別の予算を持つ}

QLoRAが向き合うのは主として既存モデルを限られた計算資源でfine-tuningする側の制約であり、FlashAttention-2はattention計算の効率、vLLM/PagedAttentionは推論時のmemory管理とservingを主眼とする。同じ「効率化」でも、最適化する段階とボトルネックは異なる。\autocite{src-1ab9d44a819a,src-5e1a3a4fc45b,src-7f830e688335,src-e03b0ba5b24e,src-fb804ac849e4}


model lifecycleを分解すると、pretrainingの巨大な一回費用だけでなく、用途別のadaptation、各requestのinference、継続的なserving capacityが別々に現れる。ある手法が一つの段階を安くしても、別の段階のcostを自動的に下げるわけではない。この分解があると、効率化技術を「速いモデル」という一語で混ぜずに比較できる。\autocite{src-1ab9d44a819a,src-5e1a3a4fc45b,src-7f830e688335,src-e03b0ba5b24e,src-fb804ac849e4}


\subsection{adaptation costを下げると用途別モデルが作りやすくなる}

QLoRAが示した方向の意味は、巨大モデルのpretrainingそのものを安価にすることではなく、既存weightsを用途に合わせて調整する際のmemory pressureを下げることにある。adaptationの入口が広がれば、単一の汎用モデルをそのまま使う以外に、domainやtaskごとの調整を検討しやすくなる。\autocite{src-7f830e688335}


ただし、memory効率が改善しても、data preparation、evaluation、deploymentまで自動化されるわけではない。fine-tuningを実行可能にする技術と、望ましいbehaviorを得るための学習設計は別問題である。2023年は「調整できること」と「調整後に何が良くなるか」を分ける必要性も強めた。\autocite{src-7f830e688335}


\subsection{同じモデルでもkernelで実行特性は変わる}

FlashAttention-2が属する層では、model architectureを変えずにattention computationの実装を改善する。これはalgorithm名やparameter数だけ見ても分からない差である。GPU上のmemory movementやparallelismの扱いが変われば、同じ数学的処理でも実行時間と利用可能なsequence長に影響しうる。\autocite{src-5e1a3a4fc45b}


この種の最適化は、モデル研究とsystems研究の境界を曖昧にする。新しいmodel capabilityを提案していなくても、既存モデルを現実的なlatencyやmemory budgetで動かせる範囲を変えるからである。2023年には「モデルの論文」と「実装の論文」を別世界として扱いにくくなった。\autocite{src-5e1a3a4fc45b}


\subsection{一件の推論と、多数の要求を捌くservingは違う}

モデルが一度応答できることと、多数の要求を安定して処理できることも別問題である。vLLM/PagedAttentionの登場は、KV-cacheを含むruntime側の資源管理が、モデル選定と同じ程度にdeployment設計へ影響しうることを可視化した。\autocite{src-1ab9d44a819a,src-e03b0ba5b24e,src-fb804ac849e4}


servingではthroughputとlatency、memory occupancy、request長、同時実行数が相互に影響する。そのためprojectが示すthroughput改善をmodel capabilityの向上として読むべきではない。runtime optimizationは「同じ能力を何人に、どの待ち時間で提供できるか」という別の性能面を扱う。\autocite{src-1ab9d44a819a,src-e03b0ba5b24e,src-fb804ac849e4}


serving layerが成熟すると、application側は巨大モデルの内部実装よりAPI surfaceを意識するようになる。同時にprovider側ではcache、batching、schedulerなどがeconomicsを左右する。tool/action layerがapplication interfaceを作ったのに対し、serving layerはそのinterfaceを継続的に成立させる裏側の基盤である。\autocite{src-1ab9d44a819a,src-e03b0ba5b24e,src-fb804ac849e4}


\subsection{小型モデルは別の制約条件を選ぶ設計}

Phi-2のような小型モデルへの関心は、parameter scaleだけを増やす以外の設計空間も示した。ここでは小型であること自体を優位性と断定せず、data qualityや用途、hardware制約と組み合わせて「必要な能力をどの資源予算で実現するか」という問いが強まったことを重視する。\autocite{src-6e61359d2fae}


フロンティア最大モデルと小型モデルは、必ずしも同じ役割を競う必要がない。cloudで最大能力を使う設計、端末や局所環境で小さなモデルを使う設計、taskに応じて複数tierを組み合わせる設計は別々に成立する。model portfolioという見方をすると、runtimeとmodel sizeは一体のdeployment decisionになる。\autocite{src-6e61359d2fae}


2023年のsystems面の変化は、モデル能力を「学習済みweightsの属性」だけで語れなくした。適応コスト、kernel、cache、serving、モデルサイズが相互に作用し、同じモデルでも運用可能性が大きく変わるという認識が、研究と実装の両側で前面に出た。\autocite{src-1ab9d44a819a,src-5e1a3a4fc45b,src-6e61359d2fae,src-7f830e688335,src-e03b0ba5b24e,src-fb804ac849e4}


\begin{claimboundary}
QLoRA、FlashAttention-2、vLLM、Phi-2に関する速度・memory・品質などの数値は、それぞれの著者やproject/vendorが設定した実験条件に依存する。異なるhardware・workload・benchmarkを横断して一つの優劣表へ変換しない。\autocite{src-1ab9d44a819a,src-5e1a3a4fc45b,src-6e61359d2fae,src-7f830e688335,src-e03b0ba5b24e,src-fb804ac849e4}
\end{claimboundary}
